\documentclass[../../../include/open-logic-section]{subfiles}

\begin{document}

\olfileid{sth}{z}{story}
\olsection{मांडणीचे आणखी तपशील}

\olref[story][approach]{sec} मध्ये आपण संच तयार होण्याच्या प्रक्रियेबद्दल शोनफिल्ड यांचे वर्णन उद्धृत केले. आता ही मांडणी थोडी अधिक नेमकी करण्यासाठी आणखी काही तत्त्वे लिहू. ती अशी:
\begin{enumerate}
\item[] \stageshier. प्रत्येक संच कोणत्या ना कोणत्या टप्प्यावर तयार होतो.
\item[] \stagesord. टप्प्यांचा क्रम असतो: काही टप्पे इतरांच्या \emph{आधी} येतात.\footnote{प्रत्यक्षात टप्पे \emph{सुक्रमित} आहेत, असे आपण अप्रकटपणे गृहीत धरू. याचा अर्थ \olref[ordinals][]{chap} मध्ये स्पष्ट केला आहे. हे महत्त्वाचे गृहीतक आहे. खरे तर \citet{Scott1974} यांच्या अतिशय कल्पक पद्धतीने हे गृहीतक प्रथम \emph{टाळता} येते आणि मग \emph{सिद्ध} करता येते. (आरंभीचा टप्पा असावा असे का मानावे, हेही त्यातून स्पष्ट होईल.) येथे त्यात जाता येणार नाही; अधिक माहितीसाठी \citet{ButtonLT1} पाहा.}
\item[] \stagesacc. कोणताही टप्पा $S$ आणि $S$ च्या \emph{आधी} तयार झालेल्या संचांचा कोणताही संग्रह घेतल्यास, टप्पा~$S$ वर नेमके तेच संच सदस्य असलेला संच तयार होतो. टप्पा~$S$ वर दुसरे काही तयार होत नाही.
\end{enumerate}
ही अनौपचारिक तत्त्वे आहेत; पण त्यांच्या आधारे झर्मेलो यांच्या संच उपपत्तीतील अनेक स्वयंसिद्धकांचे समर्थन करता येईल.

(एक सावधानतेची नोंद. पुढे आपण सर्वमान्य नावे असलेली प्रमाणित स्वयंसिद्धके मांडणार आहोत; पण नुकतीच तिरपी अक्षरे वापरून मांडलेली तत्त्वे साहित्यात कोणत्याही विशिष्ट नावांनी ओळखली जात नाहीत. आपण त्यांना केवळ उपयोगी वाटतील अशी नावे दिली आहेत.)

\end{document}
